\section{Thuật toán đặt tàu}

Giai đoạn Setup yêu cầu người chơi nhập tọa độ 6 tàu. Mỗi tàu được kiểm tra hợp lệ trước khi chấp nhận.

\subsection{Luồng đặt tàu}

\begin{enumerate}
    \item Yêu cầu người chơi nhập 4 số: \texttt{row\_bow col\_bow row\_stern col\_stern}
    
    \item Parse input thành 4 số nguyên (hỗ trợ khoảng trắng, dừng ở ký tự không hợp lệ)
    
    \item Kiểm tra hợp lệ:
    \begin{itemize}
        \item Tất cả tọa độ trong [0, 6]?
        \item Tàu nằm ngang hoặc dọc? (cùng hàng hoặc cùng cột)
        \item Độ dài = |tọa1 - tọa2| + 1 khớp với yêu cầu?
        \item Không chồng lấn (kiểm tra từng ô trên board)?
    \end{itemize}
    
    \item Nếu hợp lệ: Ghi giá trị 1 vào các ô tàu chiếm, thông báo thành công, chuyển tàu tiếp theo.
    
    \item Nếu không: Thông báo lỗi, yêu cầu nhập lại tàu này.
\end{enumerate}

\subsection{Các điều kiện hợp lệ}

\begin{table}[H]
\centering
\small
\begin{tabular}{|l|c|}
\hline
\textbf{Điều kiện} & \textbf{Kiểm tra} \\
\hline
Phạm vi tọa độ & row, col $\in$ [0, 6] \\
\hline
Hình dạng & Ngang (cùng row) hoặc Dọc (cùng col) \\
\hline
Độ dài tàu & Khớp yêu cầu (2, 3, hoặc 4) \\
\hline
Không chồng lấn & Tất cả ô chứa giá trị 0 \\
\hline
\end{tabular}
\caption{Điều kiện kiểm tra khi đặt tàu}
\end{table}

\subsection{Computer Setup (One Player mode)}

Trong chế độ One Player, Computer không yêu cầu input. Thay vào đó, board của Computer được sao chép từ một mẫu được định nghĩa sẵn (`autoBoardTemplate`) trong data segment.


\subsubsection{Công thức chuyển đổi index}

Khi người chơi cung cấp tọa độ (row, column) với row, column $\in$ [0, 6], ta tính index trong mảy 1D như sau:

\begin{equation}
\text{index} = \text{row} \times 7 + \text{column}
\end{equation}

Ví dụ:
\begin{itemize}
    \item (0, 0) $\rightarrow$ index = 0 $\times$ 7 + 0 = 0 (ô đầu tiên, góc trái trên)
    \item (0, 6) $\rightarrow$ index = 0 $\times$ 7 + 6 = 6 (ô cuối cùng hàng 0, góc phải trên)
    \item (6, 6) $\rightarrow$ index = 6 $\times$ 7 + 6 = 48 (ô cuối cùng, góc phải dưới)
    \item (3, 3) $\rightarrow$ index = 3 $\times$ 7 + 3 = 24 (ô giữa bàn cờ)
\end{itemize}

\subsubsection{Công thức tính byte offset}

Vì mỗi phần tử trong mảng là một word (4 byte), byte offset của phần tử tại index được tính:

\begin{equation}
\text{byte\_offset} = \text{index} \times 4
\end{equation}

Hoặc kết hợp cả hai:

\begin{equation}
\text{byte\_offset} = (\text{row} \times 7 + \text{column}) \times 4
\end{equation}

Trong MIPS, khi muốn đọc/ghi ô tại (row, column), ta:
\begin{enumerate}
    \item Tính index = row $\times$ 7 + column
    \item Tính byte\_offset = index $\times$ 4
    \item Dùng $t0 = $base\_address + byte\_offset để truy cập phần tử
    \item Dùng lw hoặc sw để đọc/ghi
\end{enumerate}

\subsection{Cấu trúc dữ liệu chính}

\subsubsection{Bàn cờ chính (Main Board) - boardP1, boardP2}

Lưu trạng thái tàu của mỗi người chơi. Mỗi phần tử là một 32-bit word:

\begin{itemize}
    \item \textbf{Giá trị 1}: Ô chứa phần thân tàu còn nguyên vẹn (chưa bị bắn trúng).
    \item \textbf{Giá trị 0}: Ô trống hoặc đã bị bắn trúng.
\end{itemize}

Kích thước: 49 $\times$ 4 byte = 196 byte mỗi bàn.

\subsubsection{Bản đồ bắn (Memory Board) - memP1, memP2}

Lưu lịch sử tấn công của mỗi người chơi. Mỗi phần tử là một 32-bit word:

\begin{itemize}
    \item \textbf{Giá trị 0}: Ô chưa được tấn công (unknown).
    \item \textbf{Giá trị 1}: Ô đã được tấn công nhưng trượt (MISS).
    \item \textbf{Giá trị 2}: Ô đã được tấn công và trúng (HIT).
\end{itemize}

Kích thước: 49 $\times$ 4 byte = 196 byte mỗi bản đồ.

\subsubsection{Bộ đệm input (inputBuffer)}

Một mảng ký tự (bytes) để đọc chuỗi input từ người dùng. Kích thước đủ để chứa một dòng input (khoảng 100-256 byte).

\subsubsection{Mẫu bàn cờ tự động (autoBoardTemplate)}

Một bàn cờ được xác định sẵn cho Computer trong chế độ One Player. Chứa 49 phần tử với các ô tàu đã được đặt sẵn.

\subsubsection{Biến chế độ chơi (gameMode)}

Một word 32-bit lưu chế độ chơi:
\begin{itemize}
    \item Giá trị 1: One Player vs Computer.
    \item Giá trị 2: Two Players.
\end{itemize}

\subsection{Biểu diễn trực quan trên giao diện}

Khi in bàn cờ, các ký hiệu được sử dụng để biểu diễn trạng thái ô:

\subsubsection{Bàn cờ của chính mình (Own Board)}

\begin{table}[H]
\centering
\begin{tabular}{|c|c|l|}
\hline
\textbf{Giá trị} & \textbf{Ký hiệu} & \textbf{Ý nghĩa} \\
\hline
1 & [1] & Phần tàu còn sống \\
\hline
0 (trống) & . & Ô trống, chưa bị bắn \\
\hline
0 (bị phá) & X & Ô bị bắn trúng, tàu bị phá \\
\hline
\end{tabular}
\caption{Ký hiệu trên bàn cờ của mình}
\end{table}

\textit{Chú ý}: Trên bàn cờ chính, giá trị 0 có thể biểu diễn cả ô trống và ô bị phá hủy. Để phân biệt, chương trình có thể sử dụng thêm logic hoặc sử dụng ký hiệu khác tùy thuộc vào cách in bàn cờ.

\subsubsection{Bản đồ bắn đối phương (Enemy Map)}

\begin{table}[H]
\centering
\begin{tabular}{|c|c|l|}
\hline
\textbf{Giá trị} & \textbf{Ký hiệu} & \textbf{Ý nghĩa} \\
\hline
0 & . & Chưa biết (chưa bắn) \\
\hline
1 & O & Đã bắn nhưng trượt (MISS) \\
\hline
2 & X & Đã bắn và trúng (HIT) \\
\hline
\end{tabular}
\caption{Ký hiệu trên bản đồ bắn đối phương}
\end{table}

\subsection{Ví dụ biểu diễn dữ liệu}

Giả sử người chơi 1 đặt 2 tàu:
\begin{itemize}
    \item Tàu 1 (dài 2): từ (0, 0) đến (0, 1) - nằm ngang
    \item Tàu 2 (dài 2): từ (2, 3) đến (3, 3) - nằm dọc
\end{itemize}

Bàn cờ boardP1 sẽ có giá trị 1 tại các vị trí:
\begin{itemize}
    \item Index 0 (row=0, col=0) = 1
    \item Index 1 (row=0, col=1) = 1
    \item Index 17 (row=2, col=3) = 1
    \item Index 24 (row=3, col=3) = 1
\end{itemize}

Các vị trí khác = 0.

Sau khi bắn vào (0, 0) và trúng, boardP1 tại index 0 được đặt về 0, còn memP2 (bản đồ bắn của Player 2) tại index 0 được đặt về 2 (HIT).
